\documentclass[10pt]{article}
\usepackage[letterpaper,margin=0.68in]{geometry}
\usepackage{amsmath,amssymb}
\usepackage{enumitem}
\usepackage{xcolor}
\usepackage{fancyhdr}
\usepackage{microtype}
\usepackage[most]{tcolorbox}
\definecolor{olive}{HTML}{4D5430}
\definecolor{gold}{HTML}{9A7B22}
\definecolor{pale}{HTML}{F5F1DC}
\definecolor{softgreen}{HTML}{EDF0E2}
\pagestyle{fancy}\fancyhf{}\lhead{\textcolor{olive}{MATH\& 107 Answer Key}}\rhead{\textcolor{gold}{Fall 2026}}\cfoot{\thepage}
\setlength{\headheight}{14pt}
\setlength{\parindent}{0pt}\setlength{\parskip}{5pt}
\setlist[enumerate]{leftmargin=*,itemsep=5pt,topsep=4pt}
\newtcolorbox{solbox}[1]{colback=pale,colframe=olive,boxrule=.7pt,arc=2mm,title=\textbf{#1},fonttitle=\color{olive},left=2mm,right=2mm,top=1.5mm,bottom=1.5mm}
\newcommand{\modulekey}[2]{\clearpage{\Large\bfseries\textcolor{olive}{Module #1: #2 — Answer Key}}\par\textcolor{gold}{\rule{\linewidth}{1.2pt}}\par}
\begin{document}
\modulekey{8}{Convex Geometry}
\textbf{Prequiz}
\begin{enumerate}
\item The filled rectangle is convex. A line segment between any two points in it stays in the rectangle. A crescent has pairs of points whose connecting segment leaves the set.
\item Boundary \(x+2y=8\) has intercepts \(\boxed{(8,0)}\) and \(\boxed{(0,4)}\). For \((2,3)\), \(2+2(3)=8\), so it is feasible. For \((4,3)\), \(4+6=10>8\), so it is not feasible.
\item Solve
\[
x+y=8,
\qquad 2x+y=10.
\]
Subtract to get \(x=2\), hence \(y=6\). The vertices are
\[
\boxed{(0,0),(5,0),(2,6),(0,8)}.
\]
\item
\[
P(0,0)=0,\quad P(5,0)=20,\quad P(2,6)=38,\quad P(0,8)=40.
\]
Maximum \(=\boxed{40}\) at \(\boxed{(0,8)}\).
\item
\[
f(0,0)=0,\quad f(6,0)=18,\quad f(4,3)=18,\quad f(0,5)=10.
\]
The adjacent vertices \((6,0)\) and \((4,3)\) tie, so every point on the edge joining them is optimal.
\end{enumerate}

\textbf{Matching quiz}
\begin{enumerate}
\item The filled triangle is convex; the ring is not because some connecting segments pass through the hole.
\item Boundary \(3x+y=12\) has intercepts \(\boxed{(4,0)}\) and \(\boxed{(0,12)}\). At \((3,4)\), \(3(3)+4=13>12\), so the point is not feasible.
\item Solve
\[
x+y=7,
\qquad 2x+y=10
\]
to get \((3,4)\). The vertices are
\[
\boxed{(0,0),(5,0),(3,4),(0,7)}.
\]
\item
\[
R(0,0)=0,\quad R(5,0)=25,\quad R(3,4)=31,\quad R(0,7)=28.
\]
Maximum \(=\boxed{31}\) at \(\boxed{(3,4)}\).
\item Every point on the edge joining the two adjacent optimal vertices is also optimal.
\end{enumerate}
\end{document}
